%% kapa: exam=1 round=33 year=2022 period=1 form=A subject=economics subject_ko=경제학원론 seq=34 number=74
%% source: Q-Net 감정평가사 기출문제 (cst00309) artlSeq=5209997 file=raw/1008/5209997/2022년도 제33회 감정평가사 1차시험 1교시(민법,경제학,부동산학).pdf
%% conversion: kordoc 4.24.51 -> tools/gen1.py (auto); review=proofread-sonnet; flags=none
\begin{question}{74}
소비이론에 관한 설명으로 옳은 것은?
\begin{choices}
  \item 항상소득가설(permanent income hypothesis)에 따르면, 현재소득이 일시적으로 항상소득보다 작게 되면 평균소비성향은 일시적으로 증가한다.
  \item 생애주기가설(life-cycle hypothesis)은 소비자가 저축은 할 수 있으나 차입에는 제약(borrowing constraints)이 있다고 가정한다.
  \item 케인즈 소비함수는 이자율에 대한 소비의 기간별 대체효과를 반영하고 있다.
  \item 소비에 대한 임의보행(random walk)가설은 소비자가 근시안적(myopic)으로 소비를 결정한다고 가정한다.
  \item 항상소득가설은 소비자가 차입제약에 직면한다고 가정한다.
\end{choices}
\end{question}
\answer{1}
